\documentclass[10pt]{article}
\usepackage[T1]{fontenc}
\usepackage[utf8]{inputenc}
\usepackage[french]{babel}
\usepackage{amsmath,amssymb,mathrsfs,mathpazo}
\usepackage{geometry}
\geometry{a4paper,margin=2cm}
\usepackage{graphicx,array,tabularx,tabularray,colortbl,multirow,multicol,booktabs,enumitem,pifont,ulem,fancybox,diagbox,cancel}
\usepackage{pstricks,pstricks-add,pst-plot,pst-tree,pst-node,pst-text,pst-func,pst-math,pst-bspline,pst-3dplot,pst-eucl,pst-solides3d}
\let\pscalculate\undefined
\usepackage{pst-bezier}
\usepackage{tikz,pgfplots,tkz-tab}
\usetikzlibrary{arrows,arrows.meta,patterns,calc,positioning,trees,decorations.markings,shapes.geometric,shapes.misc,shapes.arrows,intersections,decorations.pathmorphing,decorations.pathreplacing,backgrounds,shadows,babel}
\pgfplotsset{compat=1.18}
\usepgfplotslibrary{fillbetween}
\usepackage[np]{numprint}
\usepackage{listings,makecell,pas-tableur,twoopt}
\let\barre\undefined
\usepackage[pstricks]{tablvar}
\definecolor{darkgreen}{RGB}{0,100,0}
\providecommand{\red}{\color{red}}
\providecommand{\blue}{\color{blue}}
\providecommand{\green}{\color{green}}
\providecommand{\black}{\color{black}}
\providecommand{\white}{\color{white}}
\providecommand{\gray}{\color{gray}}
\providecommand{\cyan}{\color{cyan}}
\providecommand{\magenta}{\color{magenta}}
\providecommand{\yellow}{\color{yellow}}
\providecommand{\violet}{\color{violet}}
\providecommand{\orange}{\color{orange}}
\providecommand{\euro}{\texteuro}
\providecommand{\EUR}{\texteuro}
\providecommand{\ieme}{\textsuperscript{e}}
\providecommand{\degres}{\ensuremath{^\circ}}
\providecommand{\degree}{\ensuremath{^\circ}}
\providecommand{\vv}[1]{\overrightarrow{#1}}
\providecommand{\Coord}[1]{#1}
\providecommand{\decosix}{}
\providecommand{\starredbullet}{\ensuremath{\bullet}}
\pagestyle{empty}
\setlength{\parindent}{0pt}
\usepackage{adjustbox}\DeclareUnicodeCharacter{2236}{:}\DeclareUnicodeCharacter{2190}{\ensuremath{\leftarrow}}
\providecommand{\R}{\mathbb{R}}
\providecommand{\N}{\mathbb{N}}
\providecommand{\Z}{\mathbb{Z}}
\providecommand{\Q}{\mathbb{Q}}
\providecommand{\D}{\mathbb{D}}
\providecommand{\C}{\mathbb{C}}
\providecommand{\e}{{\mathrm{e}}}
\providecommand{\vect}[1]{\overrightarrow{#1}}
\providecommand{\vectt}[1]{\overrightarrow{\mathrm{#1}}}
\providecommand{\vvt}[1]{\overrightarrow{\mathrm{#1}}}
\providecommand{\np}[2][]{{#2~#1}}
\providecommand{\vv}[1]{{\overrightarrow{#1}}}
\renewcommand{\d}{{\mathrm{d}}}
\providecommand{\degres}{{^\circ}}
\providecommand{\degre}{{^\circ}}
\providecommand{\blue}{\color{blue}}
\providecommand{\red}{\color{red}}
\providecommand{\enskip}{\quad}
\providecommand{\llbracket}{\lbrack\!\lbrack}
\providecommand{\rrbracket}{\rbrack\!\rbrack}
\providecommand{\euro}{{\text{€}}}
\providecommand{\no}{{\text{n°}}}
\providecommand{\interval}[2]{{[#1;#2]}}
\providecommand{\intervFO}[2]{{[#1;#2[}}
\providecommand{\dotfill}{\ldots}

\providecommand{\ds}{\displaystyle}
\providecommand{\textding}[1]{#1}
\providecommand{\Oij}{\(\mathrm{O};\vec{\imath},\vec{\jmath}\)}
\providecommand{\Oijk}{\(\mathrm{O};\vec{\imath},\vec{\jmath},\vec{k}\)}
\providecommand{\Ouv}{\(\mathrm{O};\vec{u},\vec{v}\)}
\providecommand{\intervFF}[2]{\ensuremath{[#1;#2]}}
\renewcommand{\euro}{\ifmmode\text{\texteuro}\else\texteuro\fi}
\DeclareUnicodeCharacter{00B0}{\ifmmode^{\circ}\else\textdegree\fi}
\DeclareUnicodeCharacter{20AC}{\ifmmode\text{\texteuro}\else\texteuro\fi}
% Énoncés : Éducation nationale. Transcription et corrigés : APMEP, Denis Vergès et contributeurs. Découpage et classement : CoopMaths, Sébastien Lozano, Éric Elter et contributeurs. Adaptation de la présentation : Maths Explorer.
% CC BY-SA : https://coopmaths.fr/www/about/
% https://coopmaths.fr/alea/static/bac/2023/tex/bac_2023_03_sujet1_polynesie_3_cor.tex
\begin{document}
\noindent\textit{APMEP / CoopMaths. CC BY-SA. Présentation adaptée par Maths Explorer.}\par\medskip
\medskip

\begin{enumerate}
\item \textbf{Affirmation :} La fonction $f$ définie sur $\R$ par $f(x) = \text{e}^x - x$ est convexe.

$f$ est définie et dérivable sur $\R$ comme somme de fonctions dérivables sur $\R$.

On a $f'(x) = \text{e}^x - 1$, puis $f''(x) = \text{e}^x$.

On sait que $\text{e}^x > 0$ quel que soit $x$, donc $f$ est convexe sur $\R$. l'affirmation est vraie.

\item \textbf{Affirmation :}L'équation $\left(2\text{e}^x - 6\right)\left(\text{e}^x + 2\right) = 0$ admet $\ln (3)$ comme unique solution dans $\R$.

Cette équation se traduit par :

\[\left\{\begin{array}{l c l}
2\text{e}^x - 6&=&0\quad \text{ou}\\
\text{e}^x + 2&=&0
\end{array}\right.\]

$\bullet~~$ On sait (voir ci-dessus) que $\text{e}^x > 0 \Rightarrow \text{e}^x + 2  > 2 > 0 $. Donc la deuxième équation n'a pas de solution.

$\bullet~~$$2\text{e}^x - 6=0 \iff 2\text{e}^x =6 \iff \text{e}^x = 3 \iff x = \ln 3$, par croissance de la fonction logarithme népérien : c'est la seule solution donc l'affirmation est vraie.

\item \textbf{Affirmation :}\[\displaystyle\lim_{x \to + \infty} \dfrac{\text{e}^{2x} - 1}{\text{e}^x - x} = 0.\]

On a $\dfrac{\text{e}^{2x} - 1}{\text{e}^x - x} = \dfrac{\text{e}^{2x}\left(1 - \text{e}^{-2x} \right)}{\text{e}^{x}\left(1 - x\text{e}^{-x} \right)} = \text{e}^{x}\dfrac{\left(1 - \text{e}^{-2x} \right)}{\left(1 - x\text{e}^{-x} \right)}$.

$\bullet~~$Au numérateur on sait que $\displaystyle\lim_{x \to + \infty}\text{e}^{-2x} = 0$, donc $\displaystyle\lim_{x \to + \infty}1 - \text{e}^{-2x} = 1$ ;

$\bullet~~$Au dénominateur $x\text{e}^{-x} = \dfrac{x}{\text{e}^x}$ et l'on sait que $\displaystyle\lim_{x \to + \infty}\dfrac{x}{\text{e}^x} = 0$, donc $\displaystyle\lim_{x \to + \infty}1 - x\text{e}^{-x} = 1$.

Le quotient a pour limite 1 et comme $\displaystyle\lim_{x \to + \infty}\text{e}^{x} = + \infty$, on a finalement $\displaystyle\lim_{x \to + \infty} \dfrac{\text{e}^{2x} - 1}{\text{e}^x - x} = +\infty$ : l'affirmation est fausse.

\item La fonction $F$ est définie et dérivable sur $\R$ et sur cet intervalle :

$F'(x) = 2\text{e}^{3x} + 3(2x + 1)\text{e}^{3x} = \text{e}^{3x}(2 + 6x + 3) = (6x + 5)\text{e}^{3x} = f(x)$.

D'autre part $F(0) = 1\text{e}^{0} + 4 = 1 + 4 = 5$

$F$ est donc la primitive de $f$ sur $\R$ qui prend la valeur 5 quand $x = 0$, l'affirmation est vraie.
\item L'execution donne $\dfrac{1 + 9 + \ldots + 5}{10} = \dfrac{50}{10} = 5$.

L'affirmation est fausse.
\end{enumerate}

\bigskip

\end{document}
